\subsection{LLM-Based Structured Extraction for ESG ABSA}

\subsubsection{Prompt engineering for structured JSON output from LLMs}

Objective and design principles. The goal is to elicit highly structured, machine-parseable JSON from LLMs that encode sentence-level ESG aspect, sentiment, and provenance in a regulator-friendly schema. Prompt templates should enforce strict schema conformity, explicit field typing, and deterministic formatting to minimize downstream parsing errors. We leverage directive prompts, system roles, and constraint-enforcing prompts to bias the model toward complete, schema-compliant outputs while preserving flexibility for bilingual data \cite{10.4018/ijswis.385572,10.21203/rs.3.rs-6795706/v1,10.48550/arxiv.2510.04023} .
Template architecture. We operationalize seven prompt templates (Template 1–Template 7) that vary in prompting strategy (instruction, example-based, chain-of-thought augmentation, and retrieval-augmented prompts). Each template targets a fixed JSON schema with fields such as documentid, language, sentenceid, section, ESGpillar, aspectterm, sentimentlabel, sentimentscore, evidencespan (start/end character indices), regulatorymapping (GRI/SASB/TCFD), and provenance (source\_page, figure/table/footnote). This modular design enables ablation studies and cross-template diagnostics.
Prompt constraints and validation. Prompts embed validation hooks: after generation, a lightweight verifier checks JSON validity (JSON schema conformance), required fields presence, allowed category sets (e.g., E/S/G pillars, SASB categories), and non-negativity for confidence scores. If a mismatch occurs, the system can trigger template fallback or re-prompt with a corrected exemplar, reducing schema drift and missing fields \cite{10.18653/v1/2023.emnlp-demo.3,10.1088/2515-7620/ae1357}
Language and bilingual considerations. For Indonesian/English disclosures, prompts include language-tagging and explicit mapping cues between Indonesian terms and English equivalents, ensuring cross-language alignment within the JSON, and enabling downstream cross-language ABSA tasks. This aligns with cross-lingual ABSA literature indicating predictable gains when multilingual prompting is coupled with formal schema enforcement\cite{10.1016/j.inffus.2025.103073,10.1111/eufm.12509,10.48550/arxiv.2511.00024}

Evaluation metrics. We evaluate the prompt outputs on (i) syntactic validity (percentage of outputs that parse as valid JSON), (ii) schema completeness (fraction of required fields populated), (iii) semantic correctness (alignment of aspectterm and ESGpillar with the sentence content), and (iv) provenance fidelity (accuracy of evidencespan and regulatorymapping). A controlled corpus with ground-truth JSON annotations under both Indonesian and English is recommended for benchmarking.

\subsubsection{Zero-shot, few-shot, and chain-of-thought prompting strategies and their tradeoffs}

Zero-shot prompting. Pros: low setup cost, broad applicability to unseen documents. Cons: higher risk of hallucination, lower consistency in structured outputs, and brittle mappings to ESG taxonomies without exemplars. For ABSA, zero-shot may produce plausible yet incomplete JSON fields, requiring post hoc validation.
Few-shot prompting. Pros: improved grounding via exemplar patterns, better adherence to JSON schema, and higher accuracy in aspect term and sentiment extraction when exemplars reflect target taxonomy. Cons: risk of exemplar leakage bias and potential degradation when domain drift occurs across documents; prompts may become large and expensive.
Chain-of-thought prompting. Pros: reduces model miscomprehension by having the model expose intermediate reasoning steps leading to the final structured output, which can be truncated or logged for auditing. This can improve reliability in complex ABSA tasks where disambiguation is required (e.g., resolving whether a sentence expresses commitment vs. outcome). Cons: increased token usage and potential leakage of sensitive internal reasoning; not always essential for the final structured output and may complicate downstream parsing if not carefully trimmed.
Tradeoffs and hybrid strategies. A pragmatic strategy often combines few-shot prompting with restrained chain-of-thought for difficult sentences. We can apply chain-of-thought prompts selectively to sentences containing hedges, modalities (will, shall, intends), or quantified targets, while using direct, template-driven outputs for straightforward sentences. Empirical work on prompt strategies in ESG contexts suggests that domain-adapted prompts and few-shot exemplars improve factual extraction without incurring excessive cost or instability \cite{10.4018/ijswis.385572,10.21203/rs.3.rs-6795706/v1,10.48550/arxiv.2511.14930,10.48550/arxiv.2511.01550}.
Prompt-template selection and automation. Implement an auto-prompt manager that selects among Template 1–Template 7 based on sentence features (length, presence of numeric targets, tables vs. narrative, presence of multilingual terms). The manager logs template choice, token usage, and output quality to inform adaptive strategies in production.
\subsubsection{Known instability issues in LLM-generated structured output}

Schema drift. LLMs may progressively drift away from the exact JSON schema across long sessions or when prompts are modified, yielding outputs that are syntactically valid but semantically misaligned with the required fields (e.g., missing aspectterm, incorrect regulatorymapping). Regular schema validation and template-level regression tests are essential to detect drift early \cite{10.1037/t06541-000,10.18653/v1/2023.finnlp-2}

Missing values and partial outputs. The model may omit optional fields (evidencespan, sentimentscore) or produce placeholder values like nulls. Implementing mandatory-field enforcement within prompts and post-generation validators reduces this risk. Additionally, employing a guarded decoding strategy (e.g., strict JSON parsers with schema restoration prompts) helps recover missing fields \cite{10.18653/v1/2023.emnlp-demo.3,10.48550/arxiv.2511.14930}

Hallucination and factual drift. In ESG contexts, hallucinations may surface as fabricated regulatory mappings or spurious targets. Grounding outputs to retrieval-augmented context and enforcing external verification against standardized taxonomies (GRI/SASB/TCFD) mitigates hallucination risk. Evidence-based prompting and external knowledge lookups are recommended for high-stakes outputs \cite{10.4018/ijswis.385572,10.48550/arxiv.2511.14930,10.48550/arxiv.2511.01550}

Prompt instability across templates. Outputs can vary across different prompt formulations due to subtle wording differences. Stability-enhancing measures include using a fixed seed for any randomness, ensemble outputs across multiple prompts with majority voting, and a post-processing reconciliation step to harmonize discordant fields \cite{10.48550/arxiv.2510.04023,10.48550/arxiv.2511.14930}

\subsubsection{Multi-model comparison as a diagnostic tool}

Diagnostic paradigm. Treat model diversity as a diagnostic asset: compare a base BERT-family encoder (multilingual or Indonesian-English) against climate-focused validators (ClimateBERT with ABSA heads) and against FinBERT/ESG-BERT variants. Evaluate consistency of the JSON outputs across models for the same inputs, focusing on (i) agreement on aspectterm and ESGpillar, (ii) concordance of sentimentlabel and sentimentscore, and (iii) alignment of regulatory\_mapping.
Metrics and experiments. Use Cohen’s kappa for categorical labels (aspectterm, ESGpillar, sentimentlabel), F1 for extraction tasks, and calibration plots for sentimentscore. Conduct ablations to determine which model combination yields highest stability and regulator-friendly interpretability, and examine error concordance to identify systematic weaknesses common to certain models or languages.
Practical utility. A multi-model diagnostic workflow helps detect systematic biases, validates outputs against domain knowledge, and guides ensemble strategies (e.g., majority vote among models, or a meta-model that chooses the best-performing model per sentence type). This approach aligns with recent work on model evaluation and reliability in ESG NLP settings \cite{10.48550/arxiv.2510.22964,10.1145/3630106.3658966,10.48550/arxiv.2108.01415}.
\subsubsection{Existing ESG extraction studies using LLMs}

Summary of empirical landscape. Recent work demonstrates LLM-based approaches for ESG extraction, including QA-style information retrieval from ESG reports, retrieval-augmented generation for evidence-grounded analysis, and prompt-driven extraction of ESG indicators from long-form reports. Notable themes include domain-adaptation (FinBERT, ClimateBERT), RAG-based grounding, and the need for explainability and auditability in regulator-facing outputs. Representative examples show LLM-driven extraction of environmental commitments, governance disclosures, and risk signals, with varying degrees of grounding and verifiability \cite{10.48550/arxiv.2511.0155010.1609/aaai.v38i21.30574,10.48550/arxiv.2511.05524,10.48550/arxiv.2511.11885}
Gaps and opportunities. While LLMs enable scalable processing of ESG documents, many studies rely on single-model prompts or lack rigorous evaluation against regulatory frameworks. There is a need for standardized JSON schemas, multilingual ABSA pipelines, and regulator-aligned verification that integrates ontology-based taxonomies (GRI/SASB/TCFD) with LLM outputs. The six contribution areas outlined here aim to address these gaps by combining structured prompt engineering, retrieval grounding, and ontology-based interpretation.
Validation and reproducibility considerations. For ESG extraction work, reproducibility requires explicit documentation of prompt templates, model versions, decoding strategies, and evaluation metrics across languages and document types. The seven prompt templates and two-model configuration proposed here are designed to facilitate reproducibility and cross-linguistic comparability while enabling robust diagnostic experimentation.
\subsubsection{The seven prompt templates and two models used in this research}

Prompt templates. Template 1 through Template 7 implement a spectrum of prompting strategies (direct instruction, few-shot exemplars, chain-of-thought augmentation, region-aware prompts, regulatory-anchor prompts, evidence-tagging prompts, and explanation-enabled prompts). Each template is designed to enforce a fixed JSON schema and to inject document-region provenance, enabling precise sentence-to-JSON mappings.
Model configurations. Model A and Model B represent two distinct encoder families. Model A is a bilingual-capable transformer (e.g., XLM-Roberta or mBERT) fine-tuned on ESG-domain corpora; Model B is a climate-validated transformer (ClimateBERT-based variant) augmented with domain adapters for ESG taxonomy alignment.
Trackable outputs. For each input, outputs are parsed into the standard JSON schema, with logs capturing template ID, model used, token counts, and confidence scores. The diagnostic pipeline aggregates results across prompts and models to support stability analyses and regulator-facing traceability.
Grounding and reuse. Outputs are anchored to external ESG ontologies and standard frameworks (GRI/SASB/TCFD) via an ontology path, enabling explainability and cross-document traceability. The knowledge-graph-backed grounding provides a reproducible basis for evaluation across datasets and languages.
\subsubsection{Practical recommendations and future directions}

Secretariat-level governance for outputs. Maintain versioned prompts and models, monitor schema drift, and enforce strict JSON validation in production. Adopt a human-in-the-loop review for high-stakes extractions, especially those involving regulatory mappings.
Language-specific adaptations. Invest in bilingual ESG lexicons and term normalization pipelines to ensure Indonesian-English consistency. Integrate cross-language embeddings and align Indonesian regulatory lexicons with international ESG vocabularies within the ontology.
Evaluation protocol. Establish a standardized benchmark dataset with ground-truth JSON annotations for ESG ABSA, including both Indonesian and English samples, to enable reproducible comparisons across models and prompts.
Annotated bibliography and references block (machine-readable XML)

The references herein include works on prompt engineering, retrieval-grounded generation, RAG for ESG, and domain-adapted models in ESG contexts; you should populate the exact IEEE-formatted entries with precise DOIs and venue details in your manuscript. The candidate references listed in the Possible References section provide a spectrum of related methodologies for grounding LLM extractions in domain knowledge, reliability frameworks, and ABSA innovations.
Notes on reproducibility and ethics

All experiments should document model versions, prompt templates, decoding configurations, evaluation metrics, and languages. Where possible, share prompt templates, JSON schemas, and ground-truth corpora under appropriate licensing to enable independent replication. Given ESG data sensitivity and regulatory compliance considerations, ensure any disclosure data used in experiments respects propriety and privacy constraints.
End of section.